\scanpage{76}
le foncteur $C\longrightarrow\rho(C)$ \inserted{$\subset\underline{\mathrm{Hom}}(C,\mathrm{Ens})^\circ=\widehat{C^\circ}^{\,\circ}$} donné \deleted{par} en associant à tt objet de $C$ le foncteur covariant qu'il représente (alors il est clair que $C\longrightarrow\rho(C)$ est pl. fidèle !). On a en effet, \deleted{si $D$ est une catégorie où les $\varprojlim$ quelc. sont \unsure{repr.|requ.}, et si $\rho(C)$ est défini comme la $\Delta$-$\Delta_1$ enveloppe $\underline{\mathrm{Hom}}_{\unsure{\Delta_1|\Delta}}(C,D)\simeq\underline{\mathrm{Hom}}_{\Delta_1}(\rho(C),D)\simeq\unsure{\mathrm{H}}$}
\[
\begin{aligned}
\rho(C)&\cong T_{\rho(C),\Delta_1}(\mathrm{Ens})^\circ\\
&\cong T_{C,\Delta}(\mathrm{Ens})^\circ=\underline{\mathrm{Hom}}_\Delta(C,\mathrm{Ens})^\circ\mathrel{\therefore}\mathrm{OK}.
\end{aligned}
\]

\uline{Pb 2} $C\in\Delta$, $\Sigma\subset\mathrm{Fl}$, construire $C\Sigma^{-1}$ \uline{dans $\Delta$}, et de m\textsuperscript{me} si on se donne un ens. de couples de flèches qu'on veut égaliser. \deleted{Tout marche bien si $d'$} Soit $\rho(C,\Sigma)$ la \inserted{$\Delta_1$-}catégorie de \deleted{quotient}\inserted{fractions} de $\rho(C)$, définie par $\Sigma$, et soit $C'$ la plus petite sous-\inserted{$\Delta$-}catégorie \inserted{str. pleine} de $\rho(C,\Sigma)$ \deleted{qui} telle que $C\longrightarrow\rho(C,\Sigma)$ se factorise par $C'$, \inserted{i.e. telle que $C'$ contienne la sous-catégorie \uline{pleine} $C\Sigma^{-1}$.}

% slanted
\margin{Si $\Sigma$ contient les diagrammes des produits fibrés et des noyaux, alors il semble que $C\Sigma^{-1}$ \deleted{\unsure{pourra}}\inserted{convienne}, voir si pour tt $\delta\in d'$, les produits fibrés de type $\mathrm{Ob}\,\delta$ existent dans $C'$, i.e. \deleted{\unsure{le}} $d'$ contient le diagramme \diagram{$\begin{tikzcd}[ampersand replacement=\&,column sep=tiny,row sep=small] {} \arrow[dr] \& {} \arrow[d] \& {} \arrow[dl]\\ \& \bullet \& \end{tikzcd}$} \unsure{où l'un au moins|ou l'un analogue} \unsure{sur un ens.|est} $\times$ $\mathrm{Ob}\,\delta$. OK. Si p. ex. $d'$ est l'ens. de tous les diagrammes \uline{finis}, mais non si c'est l'ens. des diagrammes finis \uline{discrets}.}

Soit $D\in\mathrm{Ob}\,\Delta$, alors les $\Delta$-foncteurs $\varphi:C\longrightarrow D$ correspondent aux $\Delta_1$-foncteurs $\psi:\rho(C)\longrightarrow\rho(D)$ tels que $\psi$ envoie $C$ dans (la ss-cat. pleine) $D$, et les $\Delta$-$\Sigma$-foncteurs ``sont'' \deleted{tels que} les foncteurs $\psi':\rho(C,\Sigma)\longrightarrow\rho(D)$ \deleted{tels que} qui envoient $C$ dans $D$, \deleted{ou encore} ce qui revient au m\textsuperscript{me}, $C'$ dans $D$. Donc
\diagram{$\begin{tikzcd}[ampersand replacement=\&]
\underline{\mathrm{Hom}}_{\Delta,\Sigma}(C,D) \arrow[r,"\alpha"] \& \underline{\mathrm{Hom}}_\Delta(C',D) \arrow[l,bend left,"\beta"]
\end{tikzcd}$}
, et on a le foncteur restriction en sens inverse. Il est clair que les deux composés sont \deleted{égaux}\inserted{$\simeq$} les foncteurs identiques, trivial pour $\beta\alpha$, pour $\alpha\beta$ on note que \deleted{\unsure{tout d'abord} $\beta$ est pl. fidèle,} $C'$ est déduit de $C\Sigma^{-1}\underset{\mathrm{pl.}}{\subset}\rho(C,\Sigma)$ en rajoutant transfiniment des $\varprojlim$ de type $d'$, et par suite que $\underline{\mathrm{Hom}}_\Delta(C',D)\longrightarrow\underline{\mathrm{Hom}}_{\unsure{\Delta|\subset\Delta|\mathrm{Cat}}}(C,D)$ est pl. fidèle (car $\underline{\mathrm{Hom}}_\Delta(C',D)\longrightarrow\underline{\mathrm{Hom}}_{\mathrm{Cat}}(C',D)$ l'est \ldots).

\uline{NB} Dans les deux cas envisagés (Pb 1 et Pb 2) la catégorie résultante (de $I$, resp. de $C$) se déduit de la sous catégorie pleine de départ en rajoutant des $\varprojlim$ de type $d'$ (transfiniment au besoin). On trouve sauf erreur :

% slanted
\margin{il faudrait encore regarder le cas des $\unsure{\lim_\alpha|\varinjlim_\alpha|\varprojlim_\alpha}$ et du \unsure{rajoutage|v. multiple} de flèches}

\uline{Th.} Les théories algébriques de type $\Delta$ sont celles qui sont 2-représentées par des $R\in\mathrm{Ob}\,\Delta$ ayant une petite sous-catégorie \inserted{$I'$} \deleted{qui} $d'$-\inserted{co}génératrice [i.e. \deleted{tt objet de} $R$ est la réunion d'une suite transfinie de sous-catégories str. pleines de $R$ déduites en rajoutant des $\varprojlim$ de type $d'$]. Un foncteur $I\longrightarrow R$ \inserted{($I$ discret)} est basique sss il est $d'$-cogénérateur.
% archive sheet number: 75
